% !PRINTED 681
\weberpage{681}
%<<CONT

2. \gesperrt{Let $\alpha, \beta, \gamma$ be three functions in $\Omega$; then}
\weq{6}{\left(\frac{d\alpha}{d\beta}\right)\left(\frac{d\beta}{d\gamma}\right) = \left(\frac{d\alpha}{d\gamma}\right).}

For, for infinitely many points,
\[
\left(\frac{\alpha-\alpha_0}{\beta-\beta_0}\right)_0 \left(\frac{\beta-\beta_0}{\gamma-\gamma_0}\right)_0 = \left(\frac{\alpha-\alpha_0}{\gamma-\gamma_0}\right)_0,
\]
and from this Theorem 2 follows as in 1.

Accordingly, to every function $\alpha$ in $\Omega$ one can assign a function $d\alpha$ in such a way that, for any two of these functions:
\weq{7}{\left(\frac{d\alpha}{d\beta}\right) = \frac{d\alpha}{d\beta}}
becomes an actual quotient of two functions $d\alpha, d\beta$. These functions we call the \gesperrt{differentials} of $\alpha, \beta$. The differentials of the constants, and only these, are equal to zero. The other differentials are all completely determined as soon as one of them has been assumed arbitrarily (for example, constant).

3. \gesperrt{If $\alpha, \beta, \gamma, \ldots$, are arbitrary functions in $\Omega$ which satisfy an equation}
\weq{8}{F(\alpha, \beta, \gamma, \ldots) = 0}
\gesperrt{then}
\weq{9}{F'(\alpha)\,d\alpha + F'(\beta)\,d\beta + F'(\gamma)\,d\gamma + \cdots = 0.}

To prove this theorem, take a point $\mathfrak{P}$ at which none of the functions $\alpha, \beta, \gamma \ldots$ or $d\alpha, d\beta, d\gamma, \ldots$ becomes infinite or zero and at which the derived functions $F'(\alpha), F'(\beta), F'(\gamma)$ also do not vanish, and arrange, as in (3), the function $F$ according to powers and products of $\alpha - \alpha_0, \beta - \beta_0, \gamma - \gamma_0, \ldots$ It then turns out that equation (9) is satisfied for infinitely many points, and hence identically. As special cases of (9) there result:
% [sic] printed ``ordnen'' continuing ``nehme man'', a plural verb without subject; the impersonal construction requires ``ordne''
\weq{10}{\begin{aligned}
d(\alpha+\beta) &= d\alpha + d\beta,\\
d(\alpha\beta) &= \beta\,d\alpha + \alpha\,d\beta,\\
d\left(\frac{\alpha}{\beta}\right) &= \frac{\beta\,d\alpha - \alpha\,d\beta}{\beta^2}.
\end{aligned}}

The concept of the differential can be carried over without further ado to \gesperrt{holomorphic functionals}, provided the functional variables are regarded as constants. Theorem 3. and the formulae (10) which follow from it then still hold when $\alpha, \beta, \gamma, \ldots$ are holomorphic functionals and $F(\alpha, \beta, \gamma, \ldots)$ is an integral rational function of $\alpha, \beta, \gamma, \ldots$.
